\documentclass[a4paper]{article}
\usepackage[fontset=fandol]{ctex}
\usepackage{amsmath, amssymb}
\usepackage{graphicx}
\usepackage[margin=2.5cm]{geometry}
\usepackage[most]{tcolorbox}
\usepackage{etoolbox}
\usepackage{listings}
\usepackage{hyperref}
\usepackage{booktabs}
\usepackage{subcaption}
\usepackage{float}
\usepackage{tikz}
\IfFileExists{pgfplots.sty}{\usepackage{pgfplots}\pgfplotsset{compat=1.18}}{}
\newtcolorbox{knowledgebox}[1]{enhanced, colback=blue!5!white, colframe=blue!75!black, colbacktitle=blue!75!black, coltitle=white, fonttitle=\bfseries, title=#1, attach boxed title to top left={yshift=-2mm, xshift=2mm}, boxrule=1pt, sharp corners}
\newtcolorbox{importantbox}[1]{enhanced, colback=yellow!10!white, colframe=yellow!80!black, colbacktitle=yellow!80!black, coltitle=black, fonttitle=\bfseries, title=#1, sharp corners}
\newtcolorbox{warningbox}[1]{enhanced, colback=red!5!white, colframe=red!75!black, colbacktitle=red!75!black, coltitle=white, fonttitle=\bfseries, title=#1, sharp corners}
\lstset{language=Python,basicstyle=\ttfamily\small,keywordstyle=\color{blue},stringstyle=\color{red!60!black},commentstyle=\color{green!60!black},breaklines=true,frame=single,numbers=left,numberstyle=\tiny\color{gray},captionpos=b,extendedchars=false}
\newcommand{\notetitle}{Guest Lecture I}
\newcommand{\noteauthors}{Codex 基于课程页、公开视频字幕、coverage ledger 与补充官方材料整理}
\newcommand{\notedate}{\today}
\newcommand{\videochannel}{WAVLab}
\newcommand{\videopublishdate}{2023-11-27}
\newcommand{\videoduration}{unknown}
\newcommand{\videourl}{}
\newcommand{\videocoverpath}{}
\begin{document}
\begin{titlepage}
\centering
{\Large 课程笔记\par}
\vspace{1.2cm}
{\huge\bfseries \notetitle\par}
\vspace{0.8cm}
{\large \noteauthors\par}
\vspace{0.3cm}
{\large \notedate\par}
\vspace{1.2cm}
\ifdefempty{\videocoverpath}{{\small 公开视频封面不可用。\par}}{\includegraphics[width=0.82\textwidth,height=0.45\textheight,keepaspectratio]{\videocoverpath}\par}
\vfill
\begin{tcolorbox}[width=0.9\textwidth, colback=black!2!white, colframe=black!60, sharp corners]
\textbf{视频作者/频道}：\videochannel\par
\textbf{发布日期}：\videopublishdate\par
\textbf{视频时长}：\videoduration\par
\textbf{视频链接}：\href{\videourl}{\nolinkurl{\videourl}}
\end{tcolorbox}
\end{titlepage}
\tableofcontents
\newpage
\begin{warningbox}{Source Discrepancy}
2023 course page exposes only a Guest Lecture placeholder. This chapter is reconstructed from the official 2023 schedule row, the official same-course 2022 guest lecture metadata for Zhong-Qiu Wang, and the guest's official public research pages.
\end{warningbox}

\section{22.0 重建前言：源缺失与重建边界}
根据课程页主题、公开视频字幕以及本讲的 coverage ledger，22.0 重建前言：源缺失与重建边界 这一部分主要围绕 2023 course row: Guest Lecture | no public title/video/slide on the 2023 page. 展开。
写作上保持 coverage-first 原则：先把老师在这一阶段真正讲过的问题设定、模型部件、训练或推理逻辑讲清，再压缩成适合复习的结构。

\begin{figure}[H]
\centering
\includegraphics[width=0.92\textwidth]{figures/coverage_map.png}
\caption{本讲 coverage map：按 coverage\_units.jsonl 归纳出的核心主题与章节映射。}
\end{figure}

\begin{figure}[H]
\centering
\includegraphics[width=0.92\textwidth]{figures/timeline.png}
\caption{本讲 timeline：根据字幕时间范围归纳的主题推进顺序。}
\end{figure}

\subsection{2023 course row: Guest Lecture}
从证据位置看，这一部分对应的课程证据区间是 2023-11-27。课程在这里强调的核心内容可以概括为：2023 course row: Guest Lecture | no public title/video/slide on the 2023 page.。
从课程组织上看，这一部分承担的是把局部概念嵌回整条识别流水线中的作用，避免读者把公式或模块孤立地记忆。

记录系统中的附加说明指出：This chapter is an explicitly labeled reconstruction, not a verbatim transcript of a recovered 2023 public lecture.

在本章结构中，这一单元被并入“22.0 重建前言：源缺失与重建边界”，目的是把同一阶段的定义、机制和权衡集中讲清。

\subsection{本章小结}
22.0 重建前言：源缺失与重建边界 这一部分把 Guest Lecture I 中相邻的教学单元串成了一条连续叙事：既交代了这一阶段要解决的问题，也交代了方法、推理或工程权衡为什么要这样组织。

\section{22.1 Guest Lecture Reconstruction: Robust Training, Speech Enhancement, Separation, and Multichannel Processing}
根据课程页主题、公开视频字幕以及本讲的 coverage ledger，22.1 Guest Lecture Reconstruction: Robust Training, Speech Enhancement, Separation, and Multichannel Processing 这一部分主要围绕 同门课 2022 官方 guest lecture row 指出主题聚焦 robust training method、speech enhancement and separation、multichannel processing。 展开。
写作上保持 coverage-first 原则：先把老师在这一阶段真正讲过的问题设定、模型部件、训练或推理逻辑讲清，再压缩成适合复习的结构。

\subsection{同门课 2022 官方 guest lecture row 指出主题聚焦 robust training method、speech enhancement and separation、multichannel processing}
从证据位置看，这一部分对应的课程证据区间是 2022-11-28 row。课程在这里强调的核心内容可以概括为：同门课 2022 官方 guest lecture row 指出主题聚焦 robust training method、speech enhancement and separation、multichannel processing。。
从课程组织上看，这一部分承担的是把局部概念嵌回整条识别流水线中的作用，避免读者把公式或模块孤立地记忆。

记录系统中的附加说明指出：2023 course page exposes only a Guest Lecture placeholder. This chapter is reconstructed from the official 2023 schedule row, the official same-course 2022 guest lecture metadata for Zhong-Qiu Wang, and the guest's official public research pages.

在本章结构中，这一单元被并入“22.1 Guest Lecture Reconstruction: Robust Training, Speech Enhancement, Separation, and Multichannel Processing”，目的是把同一阶段的定义、机制和权衡集中讲清。

\subsection{本章小结}
22.1 Guest Lecture Reconstruction: Robust Training, Speech Enhancement, Separation, and Multichannel Processing 这一部分把 Guest Lecture I 中相邻的教学单元串成了一条连续叙事：既交代了这一阶段要解决的问题，也交代了方法、推理或工程权衡为什么要这样组织。

\section{22.2 研究主题一：Speech Enhancement and Separation}
根据课程页主题、公开视频字幕以及本讲的 coverage ledger，22.2 研究主题一：Speech Enhancement and Separation 这一部分主要围绕 讲者官方主页把研究方向明确落在 speech enhancement、speech separation、robust ASR 与 array processing。 展开。
写作上保持 coverage-first 原则：先把老师在这一阶段真正讲过的问题设定、模型部件、训练或推理逻辑讲清，再压缩成适合复习的结构。

\subsection{讲者官方主页把研究方向明确落在 speech enhancement、speech separation、robust ASR 与 array processing}
从证据位置看，这一部分对应的课程证据区间是 research summary。课程在这里强调的核心内容可以概括为：讲者官方主页把研究方向明确落在 speech enhancement、speech separation、robust ASR 与 array processing。。
从课程组织上看，这一部分承担的是把局部概念嵌回整条识别流水线中的作用，避免读者把公式或模块孤立地记忆。

记录系统中的附加说明指出：2023 course page exposes only a Guest Lecture placeholder. This chapter is reconstructed from the official 2023 schedule row, the official same-course 2022 guest lecture metadata for Zhong-Qiu Wang, and the guest's official public research pages.

在本章结构中，这一单元被并入“22.2 研究主题一：Speech Enhancement and Separation”，目的是把同一阶段的定义、机制和权衡集中讲清。

\subsection{本章小结}
22.2 研究主题一：Speech Enhancement and Separation 这一部分把 Guest Lecture I 中相邻的教学单元串成了一条连续叙事：既交代了这一阶段要解决的问题，也交代了方法、推理或工程权衡为什么要这样组织。

\section{22.3 研究主题二：Multichannel Processing}
根据课程页主题、公开视频字幕以及本讲的 coverage ledger，22.3 研究主题二：Multichannel Processing 这一部分主要围绕 TF-GridNet demo page 提供了多通道 separation / dereverberation / enhancement 的官方公开例子。 展开。
写作上保持 coverage-first 原则：先把老师在这一阶段真正讲过的问题设定、模型部件、训练或推理逻辑讲清，再压缩成适合复习的结构。

\subsection{TF-GridNet demo page 提供了多通道 separation / dereverberation / enhancement 的官方公开例子}
从证据位置看，这一部分对应的课程证据区间是 demo page。课程在这里强调的核心内容可以概括为：TF-GridNet demo page 提供了多通道 separation / dereverberation / enhancement 的官方公开例子。。
从课程组织上看，这一部分承担的是把局部概念嵌回整条识别流水线中的作用，避免读者把公式或模块孤立地记忆。

记录系统中的附加说明指出：2023 course page exposes only a Guest Lecture placeholder. This chapter is reconstructed from the official 2023 schedule row, the official same-course 2022 guest lecture metadata for Zhong-Qiu Wang, and the guest's official public research pages.

在本章结构中，这一单元被并入“22.3 研究主题二：Multichannel Processing”，目的是把同一阶段的定义、机制和权衡集中讲清。

\subsection{本章小结}
22.3 研究主题二：Multichannel Processing 这一部分把 Guest Lecture I 中相邻的教学单元串成了一条连续叙事：既交代了这一阶段要解决的问题，也交代了方法、推理或工程权衡为什么要这样组织。

\section{22.4 研究主题三：Robust Training and Data Efficiency}
根据课程页主题、公开视频字幕以及本讲的 coverage ledger，22.4 研究主题三：Robust Training and Data Efficiency 这一部分主要围绕 UNSSOR demo page 提供了 unsupervised neural speech separation 的公开研究材料，可用来支撑 robust training 相关讨论。 展开。
写作上保持 coverage-first 原则：先把老师在这一阶段真正讲过的问题设定、模型部件、训练或推理逻辑讲清，再压缩成适合复习的结构。

\subsection{UNSSOR demo page 提供了 unsupervised neural speech separation 的公开研究材料，可用来支撑 robust training 相关讨论}
从证据位置看，这一部分对应的课程证据区间是 demo page。课程在这里强调的核心内容可以概括为：UNSSOR demo page 提供了 unsupervised neural speech separation 的公开研究材料，可用来支撑 robust training 相关讨论。。
从课程组织上看，这一部分承担的是把局部概念嵌回整条识别流水线中的作用，避免读者把公式或模块孤立地记忆。

记录系统中的附加说明指出：2023 course page exposes only a Guest Lecture placeholder. This chapter is reconstructed from the official 2023 schedule row, the official same-course 2022 guest lecture metadata for Zhong-Qiu Wang, and the guest's official public research pages.

在本章结构中，这一单元被并入“22.4 研究主题三：Robust Training and Data Efficiency”，目的是把同一阶段的定义、机制和权衡集中讲清。

\subsection{本章小结}
22.4 研究主题三：Robust Training and Data Efficiency 这一部分把 Guest Lecture I 中相邻的教学单元串成了一条连续叙事：既交代了这一阶段要解决的问题，也交代了方法、推理或工程权衡为什么要这样组织。

\section{总结与延伸}
Guest Lecture I 这一讲在整门《Speech Recognition and Understanding》课程中的位置，是把 Guest Lecture I 讲成一个可复用的建模模块。
从教材化角度看，本讲最重要的不是记住单一术语，而是理解它如何嵌入完整的 automatic speech recognition pipeline：输入语音如何被表示、模型如何被训练、解码如何得到最终转写，以及这些设计分别带来什么假设和权衡。
\end{document}
